% !TeX root = ../main.tex
% Add the above to each chapter to make compiling the PDF easier in some editors.

\chapter{Implementation}\label{chapter:implementation}

The purpose of this chapter is to outline the implementational details of the library developed, the tensor permutation discovery algorithm alongside the required datastructues, and the end-to-end pipeline used for compiling the operations. We provide the reader with a top-down description of our approach, stepping down one abstraction at a time, i.e.\ from the library's interface to the core of our algorithm to how it is integrated into a state-of-the-art pipeline.

\section{Interface}

The library exposes a concise Python interface, see Figure~\ref{fig:istrategy_interface}, for permutation strategies that is independent of the concrete optimization algorithm. This interface is defined as a runtime-checkable protocol and makes it possible to register and interchange different strategy implementations without changing the surrounding compilation pipeline.

\begin{figure}[ht]
    \centering
    \begin{lstlisting}[language=python]
    @runtime_checkable
    class IStrategy(Protocol):
        @staticmethod
        def find_optimal_permutation(
            network: TensorNetwork, contraction_path: ContractionPath
        ) -> list[TensorOperation]:
            ...
    \end{lstlisting}
    \caption{Outgoing python interface for permutation discovery strategies.}
    \label{fig:istrategy_interface}
\end{figure}


This interface is the library's public contract for strategy implementations. The \texttt{find\_optimal\_permutation} method is responsible for producing an ordered sequence of tensor operations that realize an optimized contraction plan.

The strategy implementation receives two explicit inputs: a tensor network and a contraction path. The tensor network is represented by the \texttt{TensorNetwork} dataclass, see Figure~\ref{fig:tensor_network_dataclass}, which contains the initial tensor collection, the desired output indices, and a size map. The contraction path is a sequence of index pairs that specifies the pairwise contraction order. Internally, the tensor network manages its tensors through a private \texttt{\_Tensorpool} abstraction, enabling dynamic reuse of tensor objects across contraction steps instead of allocating repeated tensor storage. This decision is important for performance, and it will be highlighted again in later sections when discussing the algorithm's datastructures and execution behavior. Finally, a strategy's return type is not a raw permutation matrix or a list of index orders, but an explicit intermediate representation in the form of \texttt{TensorOperation} objects.


\begin{figure}[ht]
    \centering
    \begin{lstlisting}[language=python]
    @dataclass(init=False)
    class TensorNetwork:
        tensors: _TensorPool
        output_indices: list[int]
        size_dict: dict[int, int]
    \end{lstlisting}
    \caption{Snippet of the tensor network dataclass.}
    \label{fig:tensor_network_dataclass}
\end{figure}

The intermediate representation, see Figure~\ref{fig:intermediate_representation_objects}, decouples permutation planning from execution. It allows the optimizer to express both layout transformations and contraction steps as first-class operations that can be later compiled, scheduled, or simulated by the pipeline.

\begin{figure}[ht]
    \centering
    \begin{lstlisting}[language=python]
    class TensorOperation(ABC):
        @abstractmethod
        def apply(
            self, *inputs: TensorOperationResult, **kwargs: dict[str, Any]
        ) -> TensorOperationResult:

    class TensorPermutationOperation(TensorOperation):
        ...

    class TensorContractionOperation(TensorOperation):
        ...
    \end{lstlisting}
    \caption{Snippets of the intermeidate representation objects.}
    \label{fig:intermediate_representation_objects}
\end{figure}

Additionally, strategies expose a \texttt{get\_peak\_memory} method for a second aspect of the interface: cost estimation. It computes the maximum memory footprint that a given path and its permutations will require during execution. Similarly, \texttt{get\_total\_memory} reports the total memory movement associated with the same plan. These metrics are useful for comparing candidate strategies and selecting one that balances performance and resource usage.

By defining both permutation planning and memory estimation in the same interface, the library provides a coherent abstraction for strategy selection. Specific algorithms, such as the greedy strategy developed in this work, implement this protocol and are therefore interchangeable with alternative strategies that may be added in the future.

\section{Preparation and Peripherals}

\subsection{Persistent Contraction Path}

Before the permutation strategy can be applied, the contraction path must be prepared in a form that preserves the complete tensor network history. The library uses a dedicated \texttt{PersistentContractionPath} abstraction for this purpose. This class stores the original contraction path together with a history of tensor network states, and enforces that the history contains exactly one entry for the initial state plus one additional entry for each contraction step.

The persistent path is inspired by persistent data structures such as segment trees, where later versions share unchanged portions of earlier versions through structural sharing. In this implementation, the contraction history is built by simulating each contraction in order and appending only the changed state to the history list while preserving references to unchanged tensors wherever possible. This means that the algorithm can jump directly to any intermediate step without replaying the entire contraction sequence, and that the storage cost grows only with the number of actual changes rather than with a naive full copy of every intermediate network.

The \texttt{from\_contraction\_path} factory method is the entry point for preparing a contraction path. It takes the original tensor network and the proposed path, then constructs the sequence of states needed by the permutation strategy. The history is kept immutable from the strategy's perspective through the use of deep copies when states are retrieved, ensuring that later operations cannot accidentally mutate earlier steps. As a result, both correctness and reproducibility are preserved when the optimizer analyzes different nodes in the contraction tree.

\begin{figure}[ht]
    \centering
    \fbox{\parbox{0.8\linewidth}{Placeholder for an illustration of the persistent contraction path and state history.}} % TODO: Add Image here
    \caption{Illustration of the persistent contraction path and state history.}
    \label{fig:persistent_contraction_path}
\end{figure}

This persistent representation is therefore not only a convenience feature, but a core preparatory step for the algorithm. It is essential for the strategy's ability to inspect the tensor network at arbitrary points in the contraction sequence without recomputing the entire tree, and it ensures that the memory footprint of the preparation phase remains proportional to the number of contracted states rather than to the size of each complete tensor network copy.

\subsection{APIs and ABIs}

The library exposes a dedicated interface module under \texttt{tccg\_interface}, which acts as the bridge between the tensor contraction logic and the low-level \ac{TCCG} and runtime. At the highest level, \texttt{execute\_tccg\_contraction} orchestrates the generation of a TCCG description for a given pair of input tensors, discovers the generated function signature, compiles a \texttt{pybind11} wrapper, and finally loads and executes the resulting shared object.

This Python-side API hides the details of the \ac{TCCG} DSL, the C++ code generation pipeline, and the shared library loading semantics behind a single callable contract. The \texttt{tccg\_interface} package is implemented as a layered ABI stack, with \texttt{TCCGGenerator} responsible for producing the \texttt{.tccg} description and invoking the installed \texttt{tccg} executable, \texttt{TCCGDiscoverer} parsing the resulting generated source, \texttt{TCCGPyBind11Compiler} compiling the wrapper into a native module, and \texttt{TCCGRuntime} loading the compiled module dynamically.

Because the runtime module is created through \texttt{pybind11}, the library's ABI is defined by the exact function signature produced by \ac{TCCG} and exposed through the generated wrapper. The compilation step therefore requires access to the native \ac{HPTT} headers and libraries, which are discovered from the environment variables \texttt{HPTT\_ROOT} and \texttt{BLIS\_ROOT}. The generated \ac{TCCG} kernels are linked against \ac{HPTT} to provide the high-performance transpose and packing primitives that are necessary for the matrix-matrix multiplication kernels to achieve practical throughput.

The compiler glue is implemented in \texttt{TCCGPyBind11Compiler}, which emits a small C++ wrapper around the generated \ac{TCCG} function and then invokes the resepctive C++ compiler with the appropriate include and link paths. 

This explicit ABI binding ensures that the tensor contraction kernels generated by \ac{TCCG} are directly compatible with the \ac{HPTT}-backed transposition and workspace management used elsewhere in the library. Thus, the \texttt{tccg\_interface} package serves as the central integration point between the algorithmic layer in Python and the native tensor compute stack defined by \ac{TCCG}.
